\subsection{完备的 m-进环与逆极限}

	我们在第 6 节已经看到，若 A 为交换环，m 为 A 的理想，且 A 在 m-进拓扑下 Hausdorff 且完备，则拓扑环 A 可典范识别为离散环 \(A_{i} = A/m^{i+1} (i \in \mathbf{N})\) 关于典范映射所成的逆极限

\[
h _ {i j} \colon \mathrm{A} / \mathfrak {m} ^ {j + 1} \rightarrow \mathrm{A} / \mathfrak {m} ^ {i + 1} \quad (i \leqslant j);
\]

	注意，\(h_{ij}\) 是满射，并且若 \(n_{ij}\) 是其核，则

\[
\mathfrak {n} _ {i j} = \mathfrak {m} ^ {i + 1} / \mathfrak {m} ^ {j + 1} = (\mathfrak {m} / \mathfrak {m} ^ {j + 1}) ^ {i + 1} = (\mathfrak {n} _ {0 j}) ^ {i + 1};
\]

	特别地，\((\mathfrak{n}_{0j})^{j + 1} = 0\)。反之：

	命题 14. 设 \((\mathbf{A}_i, h_{ij})\) 是离散交换环的一个逆系统，其指标集为 \(\mathbf{N}\)，并设 \((\mathbf{M}_i, u_{ij})\) 是环逆系统 \((\mathbf{A}_i, h_{ij})\) 上模的一个逆系统。令 \(\mathfrak{n}_j\) 表示 \(h_{0j} \colon \mathbf{A}_j \to \mathbf{A}_0\) 的核，并置 \(\mathbf{A} = \varprojlim \mathbf{A}_i\)、\(\mathbf{M} = \lim \mathbf{M}_i\)。假设

\begin{enumerate}
\def\labelenumi{(\alph{enumi})}
\item
	  对所有 \(i \in \mathbf{N}\)，\(h_{i1}\) 是 \(A_i\) 上的恒等映射，并且对于 \(i \leqslant j\)，\(h_{i1}\) 和 \(u_{i1}\) 都是满射；
\item
	  对 \(i \leqslant j\)，\(h_{ij}\) 和 \(u_{ij}\) 的核分别为 \(\mathfrak{n}_j^{i+1}\) 和 \(\mathfrak{n}_j^{i+1}\mathrm{M}_j\)。
\end{enumerate}

	则：

\begin{enumerate}
\def\labelenumi{(\roman{enumi})}
\item
	  A 是完备 Hausdorff 拓扑环，M 是完备 Hausdorff 拓扑 A-模，并且典范映射 \(h_{i}: A \to A_{t}, u_{i}: M \to M_{i}\) 是满射。
\item
	  若 \(M_{0}\) 是有限生成的 \(A_{0}\)-模，则 M 是有限生成的 A-模；

	  确切地说，M 的每个有限子集 S，只要 \(u_{0}(S)\) 生成 \(M_{0}\)，就是 M 的一组生成元。
\end{enumerate}

	(i) 中的断言由《一般拓扑学》第 II 章 § 3 第 5 节命题 10 的推论和定理 1 的推论 1 得出。

	对所有 \(i \in \mathbf{N}\)，记 \(\mathfrak{m}_{i+1} = \operatorname{Ker}(h_i)\)、\(\mathbf{N}_{i+1} = \operatorname{Ker}(u_i)\)；则

\[
\mathfrak {m} _ {i + 1} = \varprojlim_ {k \geqslant 0} h _ {i, i + k} ^ {- 1} (0) = \varprojlim_ {k} \mathfrak {n} _ {i + 1} ^ {i + k}
\]

	并且 \(\mathbf{N}_{i + 1} = \lim_{k\to \infty}\mathfrak{n}_{i + k}^{i + 1}\mathbf{M}_{i + k}\)；由于 \(h_{i + k}\) 和 \(u_{i + k}\) 是满射，

\[
h _ {i + k} \left(\mathfrak {m} _ {i + 1}\right) = \mathfrak {n} _ {i + k} ^ {i + 1}, \quad u _ {i + k} \left(\mathrm{N} _ {i + 1}\right) = \mathfrak {n} _ {i + k} ^ {i + 1} \mathrm{M} _ {i + k}.\tag{16}
\]

	我们证明 \(\mathfrak{m}_i\mathrm{N}_j\subset \mathrm{N}_{i + j}\) 当 \(i\geqslant 1\)、\(j\geqslant 1\) 时成立，这等价于证明 \(u_{i + j - 1}(\mathfrak{m}_i\mathrm{N}_j) = 0\)；现在

\[
u _ {i + j - 1} (m _ {i} \mathrm{N} _ {j}) = h _ {i + j - 1} (m _ {i}) u _ {i + j - 1} (\mathrm{N} _ {j})
\]

	这等于 \(\mathfrak{n}_{i+j-1}^{i}(\mathfrak{n}_{i+j-1}^{j}\mathrm{M}_{i+j-1}) = 0\)，因为对所有 \(k \geqslant 0\)，\(\mathfrak{n}_k^{k+1}\) 是 \(h_{kk}\) 的核，且等于 0。类似地可见 \(\mathfrak{m}_i\mathfrak{m}_j \subset \mathfrak{m}_{i+j}\)。若对 \(i \leqslant 0\) 置 \(\mathfrak{m}_i = \mathrm{A}\) 且 \(\mathrm{N}_i = \mathrm{M}\)，则 \((\mathfrak{m}_i)_{i \in \mathbf{Z}}\) 是 \(\mathrm{A}\) 的滤过，\((\mathrm{N}_i)_{i \in \mathbf{Z}}\) 是 \(\mathrm{M}\) 的滤过，并且与 \(\mathrm{A}\) 上的滤过相容；\(\mathrm{A}\) 和 \(\mathrm{M}\) 上的拓扑显然就是由这些滤过定义的拓扑。于是，令 \(\mathfrak{a}\) 为 \(\mathrm{A}\) 的理想，使得 \(h_1(a) = \mathfrak{n}_1\)，并令 \(\mathrm{M}'\) 为 \(\mathrm{M}\) 中由 \(\mathrm{S}\) 生成的子模；我们将证明

\[
\mathrm{N} _ {i} = \mathfrak {a} ^ {i} \mathrm{M} ^ {\prime} + \mathrm{N} _ {i + 1} \quad \text { for } \quad i \geqslant 0.\tag{17}
\]

	记 \(a_i = h_i(a), M'_i = u_i(M')\)；只需证明

\[
u _ {i} \left(\mathrm{N} _ {i}\right) = a _ {i} ^ {t} \mathrm{M} _ {i} ^ {\prime}.
\]

	当 \(i = 0\) 时这成立，因为由假设 \(\mathbf{N}_0 = \mathbf{M}\) 且 \(\mathbf{M}_0' = \mathbf{M}_0\)。若 \(i \geqslant 1\)，则由 (16) 有 \(u_i(\mathbf{N}_i) = \mathfrak{n}_i^i\mathbf{M}_i\)。由于 \(h_{1i}\) 是满射且 \(h_{0i} = h_{01} \circ h_{1i}\)，\(h_{1i}\) 把核 \(\mathfrak{n}_i\)（即 \(h_{0i}\) 的核）映到核 \(\mathfrak{n}_1\)（即 \(h_{01}\) 的核），并且 \(\mathfrak{n}_t = \overline{h}_{1i}^{-1}(\mathfrak{n}_1)\)；于是

\[
h _ {1 i} (a _ {i}) = h _ {1} (a) = n _ {1} = h _ {1 i} (n _ {i})
\]

	又由于 \(h_{1i}\) 的核为 \(n_i^2\)，有 \(n_i \subset a_i + n_i^2\) 且 \(a_i \subset n_i\)，从而 \(n_i = a_i + n_i^2\)。此外，\(u_{0i}(M_i') = u_0(M') = M_0 = u_{0i}(M_i)\)，又由于 \(\text{Ker}(u_{0i}) = n_iM_i\)，有 \(M_t = M_t' + n_iM_t\)；从而

\[
n _ {i} ^ {t} \mathbf {M} _ {i} = (a _ {i} + n _ {i} ^ {2}) ^ {t} (\mathbf {M} _ {i} ^ {\prime} + n _ {i} \mathbf {M} _ {i}).
\]

	现在，\(\mathfrak{a}_i^k\mathfrak{n}_i^{i + 1 - k}\subset \mathfrak{n}_i^{i + 1} = 0\) 对 \(0\leqslant k\leqslant i\) 成立；于是显然 \(u_{i}(\mathbf{N}_{i}) = \mathfrak{n}_{i}\mathbf{M}_{i} = \mathfrak{a}_{i}^{i}\mathbf{M}_{t}^{\prime}\)，这就证明了 (17)。

	此外，\(m_1 = h_1(n_1)\)，从而 \(a \subset m_1\)，因此 \(a^t \subset m_1^t \subset m_i\)，从而

	\(N_{i} \subset m_{i}M' + N_{i+1}\)；另一方面显然 \(m_{i}M \subset N_{i}\)，从而 \(N_{i} = m_{i}M' + N_{i+1}\) 对所有 \(i \geqslant 0\) 成立；于是由第 9 节命题 12 的推论 3 得到 \(M' = M\)，证明完毕。

	推论 1. 在命题 14 的记号和假设下，进一步假设 \(\mathbf{M}_0\) 是有限生成的 \(\mathbf{A}_0\)-模，且理想 \(\mathfrak{n}_1\) 是 \(\mathbf{A}_1\) 的有限生成理想。令 \(\mathfrak{m}_1\) 为 \(h_0\) 的核；于是 \(\mathbf{A}\) 和 \(\mathbf{M}\) 上的拓扑分别是这个环和这个模上的 \(\mathfrak{m}_1\)-进拓扑；确切地说，对所有 \(i \geqslant 0\)，\(h_i\) 和 \(u_i\) 的核分别为 \(\mathfrak{m}_1^{i+1}\) 和 \(\mathfrak{m}_1^{i+1}\mathbf{M}\)；此外，\(\mathfrak{m}_1/\mathfrak{m}_1^2\) 是有限生成的 A-模。

	我们沿用命题 14 证明中的记号；这里的假设允许我们假定理想 a 是有限生成的。令 \(i \geqslant 0\) 为任意整数；对所有 \(j \geqslant 0\)，由 (17) 有 \(\mathbf{N}_{i+j} = \mathfrak{a}^j (\mathfrak{a}^i\mathbf{M}) + \mathbf{N}_{i+j+1} \subset \mathfrak{m}_j(\mathfrak{a}^i\mathbf{M}) + \mathbf{N}_{i+j+1}\)；反之，\(\mathfrak{m}_j(\mathfrak{a}^i\mathbf{M}) \subset \mathfrak{m}_j\mathfrak{m}_i\mathbf{M} \subset \mathfrak{m}_{i+j}\mathbf{M} \subset \mathbf{N}_{i+j}\)，从而

\[
\mathbf {N} _ {i + j} = \mathfrak {m} _ {j} (\mathfrak {a} ^ {i} \mathbf {M}) + \mathbf {N} _ {i + j + 1}.
\]

	由于 a 和 M 是有限生成的 A-模，\(a^{i}M\) 也是有限生成的 A-模。将第 9 节的命题 12 的推论 3 应用于模 \(N_{i}\) 及其滤过 \((\mathrm{N}_{ij})_{j\in\mathbf{Z}}\)，其中滤过定义为：\(N_{ij}=N_{i}\) 当 j\textless0 时；\(N_{ij}=N_{i+j}\) 当 \(j\geqslant0\) 时，由此得到 \(N_{i}=a^{i}M\)，从而 \(N_{i}\subset m_{1}^{i}M\)。但同样有 \(m_{1}^{i}M\subset m_{i}M\subset N_{i}\)，所以 \(N_{i}=m_{1}^{i}M\)。将此应用于 \(M_{i}=A_{i}\)、\(u_{ij}=h_{ij}\) 的情形，得到 \(m_{i}=m_{1}^{i}\)。此外，根据 (17)，\(m_{1}=a+m_{1}^{2}\)，这就证明了推论的最后一个断言。

	推论 2. 在推论 1 的假设下，A 为 Noether 环的充要条件是 \(A_{0}\) 为 Noether 环。

	必要性是因为 \(A_{0}\) 同构于 A 的一个商；充分性由第 10 节定理 2 的推论 5 得出。

